% MobileSec manuscript replacement blocks for the supplied main.tex
% The exact placement for each block is in OVERLEAF_MAIN_TEX_ROW_BY_ROW_GUIDE.txt.
% Remove \textcolor{blue}{...} revision marking before final submission if
% journal formatting does not require tracked changes.

% BLOCK 01
% Replace main.tex lines 237--243.
MobileSec uses a LightGBM multiclass classifier to prioritise remediation
after static analysis. The model consumes 24 scan-level attributes produced by
the cryptographic, secret-detection and network-analysis services, including
finding counts, severity distributions, vulnerability-type indicators and an
aggregate severity score. Its output is an operational remediation category,
such as weak-cipher replacement, exposed-secret removal, HTTPS enforcement or
certificate-validation repair. These categories support triage within the
platform and are mapped to relevant OWASP MASVS areas where appropriate; they
are not presented as official OWASP labels or as direct estimates of
exploitation risk.

The classification experiment uses labelled remediation outcomes retained by
the MobileSec extraction pipeline together with their provenance metadata.
Stratified training, validation and test partitions are used, and the selected
LightGBM configuration optimises validation macro-F1 to account for class-wise
performance. The reported configuration uses the multiclass objective,
\texttt{multi\_logloss} metric, 12 output classes, 7 leaves, a learning rate
of 0.08, bagging fraction of 0.8, minimum child size of 4, and regularisation
parameters $\alpha=0.3$ and $\lambda=0.6$.

% BLOCK 03
% Replace main.tex lines 476--504.
\subsection{LightGBM prioritisation evaluation}

The prioritisation model is assessed on held-out labelled remediation
outcomes. In addition to classification accuracy, macro-averaged measures are
reported because they give equal weight to each remediation class. The area
under the ROC curve (AUC) is reported using the multiclass one-versus-rest
formulation. Table~\ref{tab:lightgbm-prioritization-generated} presents the
persisted test results generated by the implemented evaluation workflow.

\begin{table}[H]
\centering
\caption{Held-out performance of the LightGBM remediation-prioritisation model.}
\label{tab:lightgbm-prioritization-generated}
\begin{tabular}{lc}
\toprule
\textbf{Metric} & \textbf{Value} \\
\midrule
Accuracy & 0.6250 \\
Macro-precision & 0.7815 \\
Macro-recall & 0.6250 \\
Macro-F1 & 0.6284 \\
AUC & 0.9468 \\
False prioritisation rate & 0.3750 \\
Under-prioritisation rate & 0.1528 \\
Over-prioritisation rate & 0.1528 \\
\bottomrule
\end{tabular}
\end{table}

These results indicate that the model captures useful ordering information for
remediation triage, while the false-prioritisation rate also shows that model
output should remain reviewable rather than being treated as an automated
security decision.

% BLOCK 04
% Insert after the Table in Block 03.
\begin{table*}[!htbp]
\centering
\caption{Examples of prioritisation errors observed in held-out predictions.}
\label{tab:lightgbm-errors}
\begin{tabular}{llllr}
\toprule
\textbf{Index} & \textbf{Expected category} & \textbf{Predicted category}
& \textbf{Interpretation} & \textbf{Priority shift} \\
\midrule
0 & \texttt{FIX\_CRYPTO\_MEDIUM} & \texttt{FIX\_WEAK\_RSA\_KEY}
& More specific cryptographic remediation predicted & +1 \\
3 & \texttt{FIX\_EXPOSED\_SECRET} & \texttt{FIX\_WEAK\_CIPHER}
& Category error at the same priority level & 0 \\
4 & \texttt{FIX\_CRYPTO\_GENERAL} & \texttt{FIX\_WEAK\_RSA\_KEY}
& More specific cryptographic remediation predicted & +1 \\
\bottomrule
\end{tabular}
\end{table*}

% BLOCK 05A
% Replace main.tex line 634.
The comparison in Table~\ref{tab:comparison_tools} is based on documented
functional scope and deployment orientation. It situates MobileSec relative to
established mobile-security platforms, but it does not by itself establish
superior detection accuracy, execution time or remediation quality. Such
claims require measurements obtained by executing each tool on the same APK
set under comparable conditions.

% BLOCK 05B
% Replace main.tex lines 636--672.
\subsection{Status of quantitative comparison with existing tools}

The present experiment provides complete measured results for MobileSec on the
DroidBench and Ghera benchmark suites. A quantitative cross-tool table is not
reported because comparable executions of MobSF, BeVigil and Yaazhini/Vooki
on the identical 283-APK set are not yet available in the reproducibility
artifacts. Once those measurements are obtained, the comparison will report
runtime, detection outcomes, precision, recall, F1-score and report-export
success using an identical evaluation protocol for each tool.

% BLOCK 10
% Replace main.tex lines 674--705, including the current benchmark table and
% its surrounding interpretation paragraphs.
\subsection{Benchmark evaluation on DroidBench and Ghera}

The benchmark experiment was conducted on 283 APKs: 118 applications from
DroidBench and 165 from Ghera. These suites provide reproducible test cases for
evaluating detector behaviour and are reported separately because their
composition and intended vulnerability scenarios differ. Archived MobileSec
scan records without reference labels are not included in the accuracy
calculation and are not used to support effectiveness claims.

\begin{table*}[!htbp]
\centering
\caption{Measured MobileSec results on the complete DroidBench and Ghera APK sets.}
\label{tab:mobilesec-benchmark-realworld-generated}
\begin{tabular}{lrr}
\toprule
\textbf{Metric} & \textbf{DroidBench} & \textbf{Ghera} \\
\midrule
Number of APKs & 118 & 165 \\
Detected findings & 484 & 901 \\
True positives & 118 & 44 \\
False positives & 0 & 121 \\
False negatives & 0 & 0 \\
Precision & 1.0000 & 0.2667 \\
Recall & 1.0000 & 1.0000 \\
F1-score & 1.0000 & 0.4211 \\
Mean scan time & 4.29 s & 8.40 s \\
Mean CPU usage & 102.4\% & 70.5\% \\
Mean RAM usage & 276.8 MB & 555.2 MB \\
Throughput & 6.71 APK/min & 6.71 APK/min \\
Concurrent APK scans & 1 & 1 \\
Mean generated export size & 77165 bytes & 80304 bytes \\
Mean report-export time & 2.28 s & 2.24 s \\
PDF report success rate & 100.0\% & 100.0\% \\
SARIF export success rate & 100.0\% & 100.0\% \\
\bottomrule
\end{tabular}
\end{table*}

The results show complete recall on both suites, with substantially lower
precision on Ghera due to false-positive detections in cases not expected to
be reported as vulnerable. This distinction is important: the benchmark
supports detection evaluation on its labelled cases, whereas operational scan
archives require independent reference labelling before accuracy measures can
be reported.

Per-category precision and recall are not reported in this version because the
benchmark annotation used for the current experiment does not yet contain a
manually reviewed category assignment for each detected finding. This analysis
is reserved for a subsequent annotation study.

% BLOCK 11
% Insert immediately before \section{Conclusions}.
\subsection{Reproducibility and developer documentation}

The benchmark workflow is reproducible from the repository services and
evaluation scripts. The reported run analyzes the complete locally supplied
DroidBench and Ghera APK collections and stores scan-detail, runtime-resource
and benchmark-label outputs used to construct
Table~\ref{tab:mobilesec-benchmark-realworld-generated}. The execution command
is:

\begin{verbatim}
node scripts/automate-paper-evaluation.js \
  --droidbench "<path-to-DroidBench-apks>" \
  --ghera "<path-to-Ghera-apks>" \
  --api "http://localhost:5000" \
  --reportgen "http://localhost:3005" \
  --concurrency 1 --derive-ground-truth
\end{verbatim}

DroidBench and Ghera remain third-party benchmark resources and must be
obtained and used according to their corresponding distribution conditions.
The repository documentation describes the service configuration, principal
endpoints and the procedure for extending the platform with additional
analysis services.

% BLOCK 14
% Replace the value in metadata row C8 at main.tex line 149.
& \url{https://github.com/ayoujil200/mobilesec/blob/main/docs/developer-guide.md} \\

% BLOCK 15
% Insert after main.tex line 459 and before the report figure.
\paragraph{Interpretation of the security score.}
The security score is designed as an interpretable triage indicator. Its
penalties encode the relative attention assigned by the platform to detected
issue types; they should not be interpreted as empirically calibrated
probabilities of compromise or as a substitute for expert risk assessment. To
examine the sensitivity of prioritisation to alternative weighting choices,
the ranking produced by the implemented weights was compared with three
alternative schemes.

\begin{table}[H]
\centering
\caption{Sensitivity of security-score ranking to alternative weighting schemes.}
\label{tab:score-sensitivity}
\begin{tabular}{lrr}
\toprule
\textbf{Weighting scheme} & \textbf{Top-10 overlap} & \textbf{Kendall's $\tau$} \\
\midrule
Implemented weights & 1.00 & 1.0000 \\
Equal weights & 0.40 & 0.7658 \\
CVSS-inspired weights & 0.90 & 0.7765 \\
Secrets-emphasised weights & 0.90 & 0.6523 \\
\bottomrule
\end{tabular}
\end{table}

Formal calibration of the score against independently assessed application
risk remains future work.

% BLOCK 17
% Insert after main.tex line 383.
\paragraph{Assessment scope of generated remediation.}
FixSuggest provides structured explanations and code-oriented remediation
proposals through the implemented user workflow. The present evaluation
demonstrates generation and presentation of these recommendations; it does
not yet measure patch correctness, developer acceptance or time-to-remediate.
These outcomes require a dedicated manual assessment protocol.

% BLOCK 19A
% In line 103 replace only the phrase:
% "supports scalable deployment and reproducible experimentation"
supports reproducible containerised deployment and benchmark-based experimentation

% BLOCK 19B
% Replace main.tex line 203.
MobileSec adopts a modular microservices architecture in which each major
stage of the analysis pipeline is implemented as a separate service. This
organization separates analysis responsibilities and supports independent
evolution of detectors and reporting components. Communication between
services combines synchronous REST-based requests with asynchronous
event-driven messaging. The architecture can support concurrent dispatch, but
the stable quantitative evaluation reported in this article uses sequential
APK execution.

% BLOCK 19C
% Replace main.tex line 474.
The platform also improves day-to-day experimental practice. Intermediate
artifacts, including manifest-derived information, certificate metadata,
permissions and extracted endpoints, are stored for reuse, which reduces
repeated preprocessing during iterative studies. The asynchronous
orchestration mechanism decouples specialised analyzers; however, the
benchmark reported here uses sequential APK processing because exploratory
parallel executions did not consistently return complete analyzer outputs.
The LightGBM prioritisation layer further moves the workflow beyond raw
finding counts toward confidence-aware triage \cite{ke2017lightgbm}.

% BLOCK 19D
% Replace main.tex line 713.
From an engineering standpoint, MobileSec demonstrates that a
microservices-based architecture can organize heterogeneous Android security
analyses while preserving reproducible deployment through containerisation.
The current implementation includes domain-specific components for
cryptography, secret exposure and network-related risk, coordinated through
persistent scan-state management and asynchronous messaging. Its operational
costs and current concurrency limitation are explicitly reported in this
article.

% BLOCK 20
% Insert after main.tex line 233 and before the Machine Learning subsection.
\paragraph{Architectural trade-offs and failure management.}
The separation of scanning, specialized detection, prioritisation and
reporting services improves modularity and makes individual components easier
to extend. It also introduces communication, deployment, monitoring and
recovery overhead that would be smaller in a single-process analyzer.

Kafka-based messaging and REST notifications require explicit treatment of
service unavailability and partial results. In the implemented workflow, the
scanner retries downstream notifications, preserves a partial or unavailable
status when a downstream service cannot provide a result, and exposes a manual
replay mechanism. A durable dead-letter queue is not part of the current
implementation.

During evaluation, exploratory execution with four APK scans in parallel
produced incomplete analyzer outcomes through time-outs or HTTP 404 responses.
Consequently, the final 283-APK benchmark was executed with one APK scan at a
time. The present study therefore does not claim a verified parallel
throughput advantage, and it does not include a measured monolithic baseline.

% BLOCK 24A
% Replace main.tex line 348.
The APK-Scanner service decompiles the application and extracts key resources,
bytecode and configuration files. It then dispatches the extracted evidence to
specialized services: CryptoCheck identifies cryptographic misuse,
SecretHunter searches for hard-coded secrets, and NetworkInspector reviews
network configurations for insecure communication patterns and TLS-related
weaknesses.

% BLOCK 26
% Insert after main.tex line 391.
Report export was also verified during the benchmark experiment described in
Table~\ref{tab:mobilesec-benchmark-realworld-generated}: PDF, JSON and SARIF
generation completed successfully for all 283 analyzed APKs.

% BLOCK 27
% Replace main.tex line 719.
The LightGBM evaluation shows that scan-level features can support
remediation-oriented triage, while the benchmark experiment provides measured
detection and reporting results for DroidBench and Ghera. Archived scans
without reference labels are treated as descriptive platform records and are
not included in accuracy claims.

% BLOCK 29
% Insert after main.tex line 722 and before Future work.
\paragraph{Limitations.}
The present evaluation has four principal limitations. First, detection
metrics are reported for DroidBench and Ghera benchmark APKs and should not be
generalized to archived or production applications without independent
reference labelling. Second, numerical comparison with MobSF, BeVigil and
Yaazhini/Vooki requires equivalent executions on the same APK set and is not
claimed in this article. Third, the security-score weighting scheme and
FixSuggest recommendations require further expert and user-centred
assessment. Finally, the reported benchmark uses sequential APK execution,
because exploratory parallel executions did not provide consistently complete
analyzer outputs.
